\begin{lemma}[Situation 7.2, Subcase 1.1 outside the stored exception]
In the setting of a
\noderef{HierarchySituationTwoCaseOneData} \(D\), write
\[
 q=C.\mathit{hierarchy}.\mathit{dividerCount},\qquad
 D_{n,q}=\noderef{HierarchyStoredColorsAt}(C.\mathit{hierarchy},q),
 \qquad
 \Gamma^q=\bigcup_{\eta\in D_{n,q}}\Gamma^q_\eta .
\]
Suppose either \(\theta\notin\Gamma^q\), or there is
\(\eta\in D_{n,q}\) such that
\(\theta\in\Gamma^q_\eta\) and \(\eta\) is not missing at
\(y_{p-1}\).  Then a contradiction follows.
\end{lemma}
\begin{proof}
Put \(H=C.\mathit{hierarchy}\),
\(A=\noderef{HierarchyStarredStage}(H,q)\),
\(\ell=\noderef{HierarchyFullTailLength}(C)\), and
\[
 B=\noderef{HierarchyFullTailPrefix}(C,\ell-2),\qquad
 B^+=\noderef{HierarchyFullTailPrefix}(C,\ell-1).
\]
Apply \noderef{HierarchyCounterexamplePenultimateElementary} to the
properness, exact one-uncoloredness, lower-rung, fewer-divider,
lexicographic-minimality, and large-set fields of
\noderef{HierarchyMinimumCounterexampleContext}.  It says exactly that
\[
                  V(B^+)\text{ is elementary under }\varphi_n.       \tag{1}
\]
The Case 1 datum makes \(\alpha\) missing at \(y_{p-1}\), whereas
\(\alpha\in D_{n,q}\).  Unfolding
\noderef{HierarchyStoredColorsAt} shows that \(\alpha\) is not missing
on \(A\).  Since \(A\preccurlyeq B\preccurlyeq B^+\) by
\noderef{TashkinovHierarchy} and
\noderef{HierarchyFullTailPrefix}, elementarity (1) makes
\(y_{p-1}\) the first displayed vertex missing \(\alpha\).

The equality
\(\Gamma^q_\alpha=\{\gamma_1,\gamma_2\}\), the avoidance field of
\noderef{GoodHierarchy}, and the preceding first-position calculation
give, for \(j=1,2\),
\[
 \gamma_j\in\overline\varphi_n(A),\qquad
 \gamma_j\notin
 \varphi_n\langle B^+-A\rangle .                              \tag{2}
\]
The tail-exclusion argument includes the divider-zero exception:
if \(A=T_n\) and the first tail edge is the exact SE/RE seed, its color
lies in the certified stored set, while each \(\gamma_j\) is
stage-missing and hence is not in \(D_{n,q}\).  Thus that individual
seed cannot be a \(\gamma_j\)-edge.  No assertion about a parallel edge
with the same endpoints is used.

Apply \noderef{ChangeLemma} to the ordered pairs
\((\gamma_j,\alpha)\).  Its Cases 1/2 alternative is supplied by the
stage-missing statement in (2), and its exact unused interval is the
second statement in (2).  All of its minimum-context hypotheses are
the fields already used above, and (1) is its derived (6.6) premise.
It follows that the full two-color components satisfy
\[
 K_{\alpha,\gamma_j}(v_{\gamma_j})
      =K_{\alpha,\gamma_j}(y_{p-1})\qquad(j=1,2),              \tag{3}
\]
where \(v_{\gamma_j}\in V(A)\) is the first
\(\gamma_j\)-missing vertex.

Invoke \noderef{HierarchySituationTwoReorderedTreesCertified} on the
underlying \noderef{HierarchySituationTwoTailData}.  Retain the exact
trees \(T^-,T^*\), their PE or exact-seed SE/RE augmentation routes,
their actual \noderef{CorrespondingExtendedTashkinovTree}
certificates, and maximum property for both replacement ETTs.  In
particular,
\[
 V(T^-)=V(T)\setminus\{y_{p-1}\},\qquad V(T^*)=V(T).          \tag{4}
\]

We next put the old hierarchy on the reordered trees.  Correspondence
identifies every positive ladder tree, raw stored set, exact connecting
set, and extension kind with the original series.  Thus the core,
closure, numerical support position, terminal missing-set equality, and
PE/non-PE clauses of \noderef{HierarchyBase} copy to bases for the two
replacement ETTs.  Copy all stages through \(T_{n,q}\), all divider
positions, and all Gamma pairs.  Replace only the terminal stage by
\(T^-\), or by \(T^*\).  The exact append relations in
\noderef{HierarchySituationTwoReorderedTrees} verify the last nested
and strict fields of \noderef{TashkinovHierarchy}; every earlier field
is unchanged.  The stored sets and stage-missing sets are unchanged,
because they use only the retained starred stages and series data.
Cardinality, pairwise disjointness, refresh, parallel-base, and
minus-closedness fields of \noderef{GoodHierarchy} therefore copy
verbatim.  Only last-stage avoidance needs checking.

If \(\theta\notin\Gamma^q\), moving the exact
\(\theta\)-edge \(e_p\) earlier crosses no Gamma color.  Removing
\((e_{p-1},y_{p-1})\) only removes possible segment colors, and in
\(T^*\) reinserting that same exact pair after \((e_p,y_p)\) creates no
new forbidden Gamma occurrence before a reserved color's first missing
vertex: (1) and the original avoidance field put every such occurrence
at or after its old cutoff.  Hence both copied hierarchies are good.

In the second alternative choose
\(\eta\in D_{n,q}\) with
\(\theta\in\Gamma^q_\eta\) and
\(\eta\notin\overline\varphi_n(y_{p-1})\).  The stored-set definition
and the avoidance field say that the first \(\eta\)-missing vertex lies
in \(B\); otherwise \(\eta\) would remain present through \(y_{p-1}\)
while its Gamma color \(\theta\) occurs on the next exact edge, contrary
to avoidance.  Thus \(e_p\) is still after the \(\eta\)-cutoff in both
reordered displays.  Pairwise disjointness of Gamma pairs shows that
this is the only pair whose avoidance could be affected by moving the
\(\theta\)-edge.  All other pairs are handled exactly as in the first
alternative.  Again both copied hierarchies are good.

We now derive the contradiction in the two possibilities for
\(\gamma_1\) at \(y_p\).  If
\(\gamma_1\in\overline\varphi_n(y_p)\), then
\(v_{\gamma_1}\in V(A)\subseteq V(T^-)\) and \(y_p\in V(T^-)\) are
distinct by vertex injectivity in \noderef{TreeSequence}.  They both
miss \(\gamma_1\), so \(V(T^-)\) is not elementary.  The corresponding
ETT, its maximum property, and the copied good hierarchy construct a
\noderef{HierarchyCounterexample} \(C^-\).  It has the same rung and
divider count as \(C\).  Unfold \noderef{HierarchyPathNumber} for its
unchanged last starred stage.  Its full tail has length \(\ell-1\), and
the last index is always a candidate, so its path number is at most
\(\ell-1\).  The datum says that the old path number is \(\ell\).
Hence the first coordinate of
\noderef{HierarchyCounterexampleMeasure} is strictly smaller, and
\noderef{HierarchyCounterexampleLess}(C^-,C)\) contradicts the
lexicographic-minimality field of the context.

It remains that \(\gamma_1\) is present at \(y_p\).  Work in \(T^*\).
By (4), its vertex set is the same as that of \(T\), so the original
non-elementarity witness, the corresponding ETT, its maximum property,
and the copied good hierarchy form a counterexample \(C^*\).  Its
full-tail length is still \(\ell\), so its path number is at most
\(\ell\).  If it is smaller, minimality is already contradicted.  We
may therefore assume it equals the old value \(\ell\), and both
coordinates of the counterexample measure then equal those of \(C\).
Transport the minimum context to \(C^*\): properness and exact
one-uncoloredness are unchanged; the corresponding-series fields
transport the complete lower-rung and large-set packages; same rung and
the copied good hierarchy transport the fewer-divider clause; and any
counterexample smaller than \(C^*\) has the same comparison with \(C\),
so the original lexicographic-minimality clause applies.

The display of \(T^*\) reaches \(y_p\) before \(y_{p-1}\); (2), and the just-checked
avoidance for the copied good hierarchy, allow a second application of
\noderef{ChangeLemma}, with the ordered pair
\((\gamma_1,\alpha)\), to this actual corresponding MP ETT and its
copied minimum data.  It gives
\[
 K_{\alpha,\gamma_1}(v_{\gamma_1})
      =K_{\alpha,\gamma_1}(y_p).                              \tag{5}
\]
Together with (3), this puts the three distinct vertices
\(v_{\gamma_1},y_{p-1},y_p\) in one two-color component, with
\(\gamma_1\) missing at the first and \(\alpha\) missing at the other
two.  By \noderef{TwoColorIncidenceDegreeAtMostTwo}, each is a
degree-at-most-one vertex in that component.  The component has a
spanning path by \noderef{TwoColorComponentSpanningPath}; a finite path
has at most two such endpoints.  Equivalently,
\noderef{TwoColorComponentsDisjoint} forbids the two terminal components
after the component through \(v_{\gamma_1}\) has been fixed.  This
contradicts (3)--(5), completing the nonexceptional branch.
\end{proof}
